% Standalone problem document, generated from the adjacent README.md.
% Compile with XeLaTeX. Mathematical content is maintained in Markdown.
\documentclass[11pt,a4paper]{article}
\usepackage[fontsize=10.5pt]{scrextend}
\usepackage[margin=22mm,headheight=15pt,headsep=9mm,footskip=11mm]{geometry}
\usepackage{fontspec}
\setmainfont{texgyrepagella}[Extension=.otf,UprightFont=*-regular,BoldFont=*-bold,ItalicFont=*-italic,BoldItalicFont=*-bolditalic]
\setsansfont{texgyreheros}[Extension=.otf,UprightFont=*-regular,BoldFont=*-bold,ItalicFont=*-italic,BoldItalicFont=*-bolditalic]
\setmonofont{texgyrecursor}[Extension=.otf,UprightFont=*-regular,BoldFont=*-bold,ItalicFont=*-italic,BoldItalicFont=*-bolditalic,Scale=MatchLowercase]
\usepackage{amsmath,amssymb}
\usepackage{unicode-math}
\setmathfont{texgyrepagella-math.otf}
\usepackage{xcolor,microtype,enumitem,needspace,fancyhdr}
\usepackage{bookmark}
\definecolor{ink}{HTML}{17344B}
\definecolor{accent}{HTML}{197D80}
\definecolor{muted}{HTML}{596773}
\hypersetup{colorlinks=true,linkcolor=accent,urlcolor=accent,pdftitle={KE-02 - Construct
a separating diagonal perturbation in nearly linear
time},pdfauthor={Open Problems in Numerical Linear Algebra}}
\urlstyle{same}
\setlength{\parindent}{0pt}
\setlength{\parskip}{5pt plus 1pt minus 1pt}
\linespread{1.04}
\setlength{\emergencystretch}{3em}
\widowpenalty=10000
\clubpenalty=10000
\displaywidowpenalty=10000
\allowdisplaybreaks[1]
\setlist{nosep,leftmargin=1.5em}
\providecommand{\tightlist}{\setlength{\itemsep}{0pt}\setlength{\parskip}{0pt}}
\providecommand{\pandocbounded}[1]{#1}
\pagestyle{fancy}
\fancyhf{}
\fancyhead[L]{\sffamily\footnotesize\color{muted}OPEN PROBLEMS IN NUMERICAL LINEAR ALGEBRA}
\fancyhead[R]{\sffamily\bfseries\footnotesize\color{accent}KE-02}
\fancyfoot[L]{\parbox[b]{0.9\textwidth}{\sffamily\fontsize{8}{10}\selectfont\color{muted}Editorial ratings; a bounded search cannot certify that a problem remains unsolved.\\Literature check: 2026-09-10}}
\fancyfoot[R]{\sffamily\footnotesize\color{muted}\thepage}
\renewcommand{\headrulewidth}{0.3pt}
\renewcommand{\headrule}{\hbox to\headwidth{\color{accent}\leaders\hrule height \headrulewidth\hfill}}
\makeatletter
\renewcommand{\section}{\Needspace{5\baselineskip}\@startsection{section}{1}{0pt}{12pt plus 2pt minus 2pt}{4pt}{\sffamily\bfseries\large\color{ink}}}
\renewcommand{\subsection}{\Needspace{4\baselineskip}\@startsection{subsection}{2}{0pt}{9pt plus 1pt minus 1pt}{3pt}{\sffamily\bfseries\normalsize\color{ink}}}
\makeatother
\setcounter{secnumdepth}{0}
\begin{document}
{\sffamily\small\color{muted}Eigenvalues and inverse problems\par}
\vspace{3pt}
{\raggedright\hyphenpenalty=10000\exhyphenpenalty=10000\sffamily\bfseries\fontsize{21}{24}\selectfont\color{ink}Construct
a separating diagonal perturbation in nearly linear time\par}
\vspace{7pt}
{\sffamily\small\textbf{Difficulty:} challenging\quad\textbf{Importance:} interesting
to the community\par}
{\sffamily\small\bfseries\color{accent}Status: Open\par}
\vspace{4pt}
{\color{accent}\hrule height 0.8pt}
\vspace{6pt}
\textbf{Topic:} Hermitian eigenproblems; deterministic regularization

\textbf{Rating rationale:} Challenging reflects derandomizing eigenvalue
separation within a nearly linear cost; community impact is
deterministic preprocessing for Hermitian eigensolvers.

\subsection{Context and notation}\label{context-and-notation}

All norms are Euclidean vector norms or their induced matrix norms.

\subsection{Problem statement}\label{problem-statement}

For a Hermitian matrix \(H\) with eigenvalues listed with multiplicity,
define

\[
\operatorname{gap}(H)=\min_{i\ne j}|\lambda_i(H)-\lambda_j(H)|.
\]

Do universal constants \(a,c_0,C,q>0\) and a deterministic algorithm
exist with the following property? For every \(n\geq2\), Hermitian
tridiagonal \(T\in\mathbb C^{n\times n}\) with \(\|T\|_2\leq1\), and
\(\delta\in(0,1/2)\), the algorithm receives the three diagonals and
returns a real diagonal matrix \(D\) such that

\[
\|D\|_2\leq\delta,
\qquad
\operatorname{gap}(T+D)\geq c_0(\delta/n)^a,
\]

using at most \(Cn[1+\log(n/\delta)]^q\) exact arithmetic operations and
comparisons? Square roots of nonnegative real numbers are allowed at
unit cost; extracting bits from exact real inputs is not.

This selects the explicitly proposed nearly linear version of the
deterministic Minami problem. It preserves tridiagonal structure and
requires only the perturbation as output. Its diagonal restriction and
cost target make it distinct from the general, cubic-cost regularization
question
\href{https://github.com/ajt60gaibb/OpenProblemsInNLA/blob/main/eigenvalues-and-inverse-problems/IE-07/README.md}{IE-07}.

\subsection{References}\label{references}

Amsel et al.,
\href{https://arxiv.org/html/2602.05394v3\#S3.SS1}{workshop report},
Problem 3.2. Sobczyk,
\href{https://arxiv.org/html/2410.21550v2}{\emph{Deterministic
complexity analysis of Hermitian eigenproblems}}, §1.1 for the
arithmetic model, Theorem 1.1 for the tridiagonal eigensolver, and §5,
question 4, for deterministic spectral separation.

\subsection{Earlier status check ---
2026-09-08}\label{earlier-status-check-2026-09-08}

Searches for \textit{deterministic Minami perturbation},
\textit{Minami deterministic eigenvalue gap algorithm 2026}, and
\textit{Minami derandomization matrix} found no construction meeting
this bound. Sobczyk's April 2025 revision gives a deterministic full
tridiagonal diagonalization in nearly quadratic time; it does not
provide the requested nearly linear diagonal perturbation. The workshop
retains Problem 3.2 in its August 2026 update. The normalization and a
bounded perturbation are stated explicitly here, making the scale of its
gap guarantee unambiguous.

\subsection{Audit update --- 2026-09-10}\label{audit-update-2026-09-10}

Rechecked \href{https://arxiv.org/html/2602.05394v3}{workshop version
3}, Problem 3.2. Searches for deterministic Minami-type separation and
nearly linear tridiagonal perturbation constructions found no solution
of the simultaneous separation and runtime requirements.
\end{document}
